\chapter{随机变量的数字特征}

之前讨论的分布律和密度函数是随机变量最完整的刻画，但很多情况下，分布规律是无法得知的，或者我们只能根据经验猜测可能符合某个分布规律。
在这种情况下，我们退而求其次，考察它的均值和离散程度等一些指标，这对生产生活也能提供指导意义。
这些特征在概率论中称为“数字特征”。

本章要点：
\begin{itemize}
    \item 数学期望。
    \item 方差。
    \item 协方差。
    \item 相关系数。
\end{itemize}

\newpage
\section{数学期望和中位数}

本节介绍一维随机变量的第一个重要数字特征——数学期望，也即平均值，顺带介绍中位数。

本节要点：
\begin{itemize}
    \item 掌握数学期望的概念；
    \item 理解数学期望的数学意义和物理意义；
    \item 了解中位数的概念。
\end{itemize}

%============================================================
\subsection{一维随机变量数学期望的概念}

\begin{definition}[一维随机变量的数学期望]
设离散变量$X$有分布律$p_i$，若无穷级数$\sum_{i=1}^{+\infty}{x_i\cdot p_i}$绝对收敛（即$\sum_{i=1}^{+\infty}{\left| x_i\cdot p_i \right|}$收敛），则称该无穷级数为$X$的{\bf 数学期望}或{\bf 均值}，记作$E\left( X \right) $或$EX$。
设连续变量$X$有密度函数$f\left( x \right) $，若反常积分$\int_{-\infty}^{+\infty}{xf\left( x \right) \cdot dx}$绝对收敛（形式同上），则称该积分为$X$的{\bf 数学期望}或{\bf 均值}，同样记$E\left( X \right) $或$EX$：
\begin{align*}
&EX:=\sum_{i=1}^{+\infty}{x_i\cdot p_i} \\
&EX:=\int_{-\infty}^{+\infty}{xf\left( x \right) \cdot dx}
\end{align*}
若离散变量取值为有限个，则必有数学期望，若连续变量为定积分，也必有数学期望。
\end{definition}

一般，我们用$E\left( X \right) $表示谓词，即计算$X$的数学期望这个过程，用$EX$表示名词，即数学期望这个值。
$E$是Expectation的首字母。

\begin{tcolorbox}
级数有个现象，如果调换各项的顺序，收敛性可能发生改变，即便收敛，也可能收敛到不同的值，而在概率论中，随机变量的取值和顺序无关，所以这里要求绝对收敛。
\end{tcolorbox}

\begin{definition}[一维随机变量函数的数学期望]
设一维离散变量$X$有分布律$p_i$，若有另一随机变量$Y=g\left( X \right) $，且无穷级数$\sum_{i=1}^{+\infty}{g\left( x_i \right) p_i}$绝对收敛，则称$E\left[ g\left( X \right) \right] $为{\bf 函数$Y=g\left( X \right) $的数学期望}。
同样定义，设连续变量$X$有密度函数$f\left( x \right) $，若有另一随机变量$Y=g\left( X \right) $，且反常积分$\int_{-\infty}^{+\infty}{g\left( x \right) f\left( x \right) \cdot dx}$绝对收敛，则称$E\left[ g\left( X \right) \right] $为{\bf 函数$Y=g\left( X \right) $的数学期望}：
\begin{align*}
&EY:=\sum_{i=1}^{+\infty}{g\left( x_i \right) \cdot p_i} \\
&EY:=\int_{-\infty}^{+\infty}{g\left( x \right) f\left( x \right) \cdot dx}
\end{align*}
\end{definition}

该公式的意义在于，有时候我们很难获得$Y$的具体分布规律，但可以通过已知的$X$和$Y=g\left( X \right) $获得$Y$的数学期望。

\begin{theorem}
设一维随机变量$X$有分布函数$F\left( x \right) $和密度函数$f\left( x \right) $，若$f\left( x \right) $连续，则$E\left[ F\left( x \right) \right] =\frac{1}{2}$。
\end{theorem}

\begin{proof}
\begin{align*}
    E\left[ F\left( x \right) \right] &=\int_{-\infty}^{+\infty}{F\left( x \right) f\left( x \right) \cdot dx}=\int_{-\infty}^{+\infty}{F\left( x \right) \cdot F'\left( x \right) \cdot dx}\\
    &=\int_{-\infty}^{+\infty}{F\left( x \right) \cdot d\left( F\left( x \right) \right)}=\left. \frac{1}{2}F^2\left( x \right) \right|_{-\infty}^{+\infty}=\frac{1}{2}
\end{align*}
\end{proof}

%============================================================
\subsection{二维随机变量数学期望的概念}

\begin{definition}[二维随机变量的数学期望]
设二维离散变量$\left( X,Y \right) $有分布律$p_{ij}$，若下列式子中相应的无穷级数绝对收敛，则称该无穷级数为{\bf 二维离散变量对$X$（或$Y$）的数学期望}，即：
\begin{align*}
&EX:=\sum_i{\sum_j{x_i\cdot p_{ij}}}=\sum_i{\left( x_i\cdot \sum_j{p_{ij}} \right)}=\sum_i{x_i\cdot p_{i\cdot}} \\
&EY:=\sum_j{\sum_i{y_j\cdot p_{ij}}}=\sum_j{\left( y_j\cdot \sum_i{p_{ij}} \right)}=\sum_j{y_j\cdot p_{\cdot j}}
\end{align*}
同样定义，设二维连续变量$\left( X,Y \right) $有密度函数$f\left( x,y \right) $，若下列式子中相应的反常积分绝对收敛，则称该积分为{\bf 二维连续变量对$X$（或$Y$）的数学期望}，即：
\begin{align*}
EX:&=\iint\limits_{\mathbb{R} \times \mathbb{R}}{xf\left( x,y \right) \cdot d\sigma}=\int_{-\infty}^{+\infty}{\left( xdx\int_{-\infty}^{+\infty}{f\left( x,y \right) \cdot dy} \right)} \\
&=\int_{-\infty}^{+\infty}{xf_X\left( x \right) \cdot dx} \\
EY:&=\iint\limits_{\mathbb{R} \times \mathbb{R}}{yf\left( x,y \right) \cdot d\sigma}=\int_{-\infty}^{+\infty}{\left( ydy\int_{-\infty}^{+\infty}{f\left( x,y \right) \cdot dx} \right)} \\
&=\int_{-\infty}^{+\infty}{yf_Y\left( y \right) \cdot dy} \\
d\sigma &=dxdy
\end{align*}
\end{definition}

\begin{tcolorbox}
对比一维变量的数学期望：
\begin{align*}
&EX:=\sum_{i=1}^{+\infty}{x_i\cdot p_i} \\
&EX:=\int_{-\infty}^{+\infty}{xf\left( x \right) \cdot dx}
\end{align*}
二维变量对某个变量的数学期望和一维变量在原理上一致，只是分布从一维变量的分布变成了二维变量的边缘分布。
\end{tcolorbox}

\begin{definition}[二维随机变量函数的数学期望]
若二维离散变量$\left( X,Y \right) $有分布律$p_{ij}$，若$Z=z\left( X,Y \right) $且$\sum_{i=1}^{+\infty}{\sum_{j=1}^{+\infty}{z\left( x_i,y_j \right) \cdot p_{ij}}}$绝对收敛，则称该级数为{\bf 二维变量函数$Z=z\left( X,Y \right) $的数学期望}，记作$E\left( Z \right) $或$EZ$。
同样定义，设二维连续变量$\left( X,Y \right) $有密度函数$f\left( x,y \right) $，若$Z=z\left( X,Y \right) $且反常积分$\iint\limits_{\mathbb{R} \times \mathbb{R}}{z\left( x,y \right) f\left( x,y \right) \cdot d\sigma},  d\sigma =dxdy$绝对收敛，则称该积分为{\bf 二维变量函数$Z=z\left( X,Y \right) $的数学期望}，同样记作$E\left( Z \right) $或$EZ$：
\begin{align*}
&EZ:=\sum_i{\sum_j{z\left( x_i,y_j \right) \cdot p_{ij}}} \\
&EZ:=\iint\limits_{\mathbb{R} \times \mathbb{R}}{z\left( x,y \right) f\left( x,y \right) \cdot d\sigma} \\
&d\sigma =dxdy
\end{align*}
\end{definition}

在二维变量中，我们很少考虑单变量的期望，而是考虑两个变量关系的期望，如$X+Y,XY$等。

%============================================================
\subsection{数学期望的数学意义}

数学期望是随机变量的加权和，权重是概率，由于总概率为1，于是$EX<x_{+\infty}$，所以期望也称加权平均。
任何一个给定的随机变量，若有数学期望，则都是一个固定的、可计算的常量。

从线性代数看，数学期望是随机变量构成的向量及其分布规律构成的向量的标量积。以离散型随机变量为例：
\begin{align*}
&EX:=\sum_{i=1}^{+\infty}{x_ip_i}=\boldsymbol{x}^T\boldsymbol{p} \\
&\begin{cases}
	\boldsymbol{x}=\left( x_1,x_2,\cdots ,x_n,\cdots \right)\\
	\boldsymbol{p}=\left( p_1,p_2,\cdots ,p_n,\cdots \right)\\
\end{cases}
\end{align*}
$EX$越大表示$\boldsymbol{p}$和$\boldsymbol{x}$的契合度越高，或者说越靠后的事件越容易发生。

\begin{tcolorbox}
这里的数学意义更多是为意义而意义，是一个“形象的比喻”，不用太过深究。
\end{tcolorbox}

%============================================================
\subsection{数学期望的物理意义}

先考虑一维物体的转动，如下图。长度为$L$的细长杆质量$M$且分布不均匀，有线密度$\rho \left( x \right) ,x\in \left[ d,d+L \right] $，一端距原点$d$，绕原点匀速转动且有角速度$\omega $，分析所需的向心力。可以认为是同质量的质点在$x_0$处以通常角速度转动所需的向心力。
\begin{figure}[h]
\centering
\includegraphics[height=3cm]{4.1.png}
\end{figure}

可以假设细长杆的一个微元可视为质点，则向心力微元和整体向心力：
\begin{align*}
&dF=\rho \left( x \right) dx\cdot \omega ^2\cdot x\quad \left[ \mathrm{N} \right] \\
&F=\int_d^{d+L}{dF}=\int_d^{d+L}{\rho \left( x \right) dx\cdot \omega ^2\cdot x}=\omega ^2\int_d^{d+L}{x\rho \left( x \right) \cdot dx}\quad \left[ \mathrm{N} \right]
\end{align*}
若将其视为质点，则向心力等价为质量$M$的质点小球在$x_0$处以角速度$\omega $匀速转动，有：
\[
\omega ^2\int_d^{d+L}{x\rho \left( x \right) \cdot dx}=M\cdot \omega ^2\cdot x_0\quad \left[ \mathrm{N} \right]
\]
于是可得：
\[
x_0=\frac{\int_d^{d+L}{x\rho \left( x \right) \cdot dx}}{M}=\int_d^{d+L}{x\frac{\rho \left( x \right)}{M}\cdot dx}\quad \left[ \mathrm{m} \right]
\]
对比一维变量的数学期望定义式：
\[
EX=\int_{-\infty}^{+\infty}{xf\left( x \right) \cdot dx}
\]
{\bf 数学期望的物理意义为绕原点匀速转动时的质心位置，有量纲$\left[ \mathrm{m} \right] $}。

简单来讲就是细长杆的质心，之所以要特别强调转动，是为了下面的二维考虑，其次，在考察方差的物理意义时，还会联系转动惯量和平行轴定理。

考虑二维平面，如下图，面密度$\rho \left( x,y \right) $的平面$D$，绕{\it y}轴匀速转动$\omega $ ，在{\it x}轴方向的向心力$F$等效到质点小球$M$产生的向心力。
\begin{figure}[h]
\centering
\includegraphics[height=3cm]{4.2.png}
\end{figure}
\begin{align*}
F&=\iint\limits_D{dF}=\iint\limits_D{\rho \left( x,y \right) d\sigma \cdot \omega ^2\cdot x} \\
&=\omega ^2\iint\limits_D{x\rho \left( x,y \right) \cdot d\sigma}=M\cdot \omega ^2\cdot x_0 \quad \left[ \mathrm{N} \right]
\end{align*}
对比二维变量的数学期望定义式：
\[
EX=\iint\limits_{\mathbb{R} \times \mathbb{R}}{xf\left( x,y \right) \cdot d\sigma}
\]
{\bf 二维变量$\left( X,Y \right) $的数学期望$EX$的物理意义为平面绕{\it y}轴匀速转动时在{\it x}方向的质心位置，有量纲$\left[ \mathrm{m} \right] $。}
$EY$的物理意义同样可得，略。

\begin{tcolorbox}
对于函数的数学期望，由于函数无固定形态，所以量纲也无具体表达。
\end{tcolorbox}

%============================================================
\subsection{数学期望的性质}

{\bf 性质1}：常数的数学期望为其本身，即$E\left( c \right) =c$。

{\bf 性质2}：$X\geqslant a\Rightarrow EX\geqslant a,X\leqslant a\Rightarrow EX\leqslant a$

{\bf 性质3}：$E\left( kX \right) =k\cdot E\left( X \right) $

{\bf 性质4}：$E\left( kX+c \right) =k\cdot E\left( X \right) +c$，由性质1+3可证。

{\bf 性质5}：设$\left( X,Y \right) $为二维随机变量，则$E\left( X\pm Y \right) =E\left( X \right) \pm E\left( Y \right) $

{\bf 性质6}：设$\left( X,Y \right) $为二维随机变量，若$X,Y$相互独立$\Rightarrow E\left( XY \right) =E\left( X \right) \cdot E\left( Y \right) $

{\bf 性质7}：设$\left( X,Y \right) $为二维随机变量，则$\left[ E\left( XY \right) \right] ^2\leqslant E\left( X^2 \right) \cdot E\left( Y^2 \right) $

~

性质1表示了常数$c$的分布，其实就是质点小球，自然质心位置为$c$。
\begin{figure}[h]
    \begin{minipage}{0.48\linewidth}
    \centerline{\includegraphics[width=4.5cm]{4.3.png}}
    \end{minipage}
\hfill
    \begin{minipage}{0.48\linewidth}
    \[
    X\sim \left( \begin{array}{l} c\\ 1\\ \end{array} \right)
    \]
    \end{minipage}
\end{figure}

性质2表示了质心的范围，物理意义也很明显，质心位置肯定不会落到物体之外。
\begin{figure}[h]
\centering
\includegraphics[height=2.5cm]{4.4.png}
\end{figure}

性质3表达当{\it x}轴拉伸时，质心也有同样的拉伸，压缩亦然。
\begin{figure}[h]
\centering
\includegraphics[height=2.5cm]{4.5.png}
\end{figure}

\begin{tcolorbox}
这三个性质可结合物理意义理解。
\end{tcolorbox}

性质5用定义即可证明：
\begin{align*}
E\left( X\pm Y \right) &=\sum_{i=1}^{+\infty}{\sum_{j=1}^{+\infty}{\left( x_i\pm y_j \right) \cdot p_{ij}}}=\sum_{i=1}^{+\infty}{\sum_{j=1}^{+\infty}{\left( x_ip_{ij}\pm y_jp_{ij} \right)}} \\
&=\sum_{i=1}^{+\infty}{\sum_{j=1}^{+\infty}{x_ip_{ij}}}\pm \sum_{i=1}^{+\infty}{\sum_{j=1}^{+\infty}{y_jp_{ij}}} \\
&=E\left( X \right) \pm E\left( Y \right)
\end{align*}

性质6，独立则有$p_{ij}=p_{i\cdot}\cdot p_{\cdot j}$，于是：
\begin{align*}
E\left( XY \right) &=\sum_{i=1}^{+\infty}{\sum_{j=1}^{+\infty}{\left( x_iy_j \right) \cdot p_{ij}}}=\sum_{i=1}^{+\infty}{\sum_{j=1}^{+\infty}{\left( x_iy_j \right) \cdot p_{i\cdot}\cdot p_{\cdot j}}} \\
&=\sum_{i=1}^{+\infty}{\left( x_ip_{i\cdot}\cdot \sum_{j=1}^{+\infty}{y_jp_{\cdot j}} \right)}=\sum_{i=1}^{+\infty}{\left[ x_ip_{i\cdot}\cdot E\left( Y \right) \right]} \\
&=E\left( X \right) \cdot E\left( Y \right)
\end{align*}

\begin{tcolorbox}
数学期望的性质常用于随机变量组合成新变量的数学期望的计算。
由于数学期望是线性运算，所以有
\[
E\left( aX+bY+c \right) =aEX+bEY+c
\]
但对于非线性组合，会有一些前置条件或者为不等式。
\end{tcolorbox}

\begin{tcolorbox}
$\left( EX \right) ^2$不大于$E\left( X^2 \right) $，证明关系留到介绍方差后再给出。
\end{tcolorbox}

%============================================================
\subsection{数学期望的范围}

数学期望的范围是变量的范围，这点从数学期望的物理意义也很好理解，特别是当变量的分布为常数分布时，取到最大值。
此时，随机事件变成确定性事件。对应地，函数的数学期望的范围是函数的值域。
\begin{figure}[h]
    \begin{minipage}{0.48\linewidth}
    \centerline{\includegraphics[width=4.5cm]{4.3.png}}
    \end{minipage}
\hfill
    \begin{minipage}{0.48\linewidth}
    \[
    X\sim \left( \begin{array}{l} c\\ 1\\ \end{array} \right)
    \]
    \end{minipage}
\end{figure}

%============================================================
\subsection{中位数的概念}

\begin{definition}[中位数]
若随机变量$X$有分布函数$F\left( x \right) $，若$m$满足
\[
F\left( m \right) =\frac{1}{2}
\]
则称$m$为{\bf 随机变量$X$的中位数}。
\end{definition}

数学期望和中位数都反映了随机变量的一个均值情况。中位数表示了该点前后的概率是一样的，优点是受特大值、特小值影响小，而且必然存在，缺点是处理不方便（如两个随机变量加减乘除后的中位数），而且不唯一，特别在离散型随机变量，由于取值的离散型，中位数大多不会精确的存在。

%============================================================
\subsection{例}

\begin{example}
若盒子中有2个红球1个白球，先分别采取不放回抽取和放回抽取两种策略，记$X$为首次抽中红球时的取球的次数，分别计算$EX$。
\end{example}

解：

1. 不放回抽取时，$X$只能为$1,2$，且易得$P\left\{ X=1 \right\} =\frac{2}{3},P\left\{ X=2 \right\} =\frac{1}{3}\cdot \frac{2}{2}=\frac{1}{3}$，于是：
\[
EX=1\cdot \frac{2}{3}+2\cdot \frac{1}{3}=\frac{4}{3}
\]

2. 放回抽取时，$X=1,2,\cdots $，且有$P\left\{ X=i \right\} =\left( \frac{1}{3} \right) ^{i-1}\cdot \frac{2}{3}$，于是$EX=\frac{2}{3}\cdot \sum_{i=1}^{+\infty}{i\cdot \left( \frac{1}{3} \right) ^{i-1}}$，对于该无穷级数，等号两边都乘以$1/3$：
\begin{align*}
&\because \frac{1}{3}EX=\frac{2}{3}\cdot \sum_{i=1}^{+\infty}{i\cdot \left( \frac{1}{3} \right) ^i} \\
&\begin{aligned}
	\therefore \frac{2}{3}EX&=EX-\frac{1}{3}EX=\frac{2}{3}\cdot \left[ \sum_{i=1}^{+\infty}{i\cdot \left( \frac{1}{3} \right) ^{i-1}}-\sum_{i=1}^{+\infty}{i\cdot \left( \frac{1}{3} \right) ^i} \right]\\
	&=\frac{2}{3}\cdot \sum_{i=1}^{+\infty}{\left( \frac{1}{3} \right) ^{i-1}}=\frac{2}{3}\cdot \frac{1}{1-\frac{1}{3}}=1\\
\end{aligned} \\
&\therefore EX=\frac{3}{2}
\end{align*}

~

\begin{example}
若变量$X$的密度函数$f\left( x \right) =2x,x\in \left( 0,1 \right) $，求$X$和$X^2$的数学期望。
\end{example}

解：
\begin{align*}
&E\left( X \right) =\int_{-\infty}^{+\infty}{xf\left( x \right) \cdot dx}=\int_0^1{x\cdot 2x\cdot dx}=\left. \frac{2}{3}x^3 \right|_{0}^{1}=\frac{2}{3} \\
&E\left( X^2 \right) =\int_{-\infty}^{+\infty}{x^2f\left( x \right) \cdot dx}=\int_0^1{x^2\cdot 2x\cdot dx}=\left. \frac{1}{2}x^4 \right|_{0}^{1}=\frac{1}{2}
\end{align*}

~

\begin{example}
设离散随机变量$X$和连续随机变量$Y$如下
\begin{align*}
&X~\left( \begin{matrix}
    0&		10\\
    \frac{3}{5}&		\frac{2}{5}\\
\end{matrix} \right) \\
&f\left( y \right) =\frac{1}{2}e^{-\left| y \right|}
\end{align*}
且$X,Y$相互独立，求$E\left( XY^2-2X^2Y+1 \right) $。
\end{example}

解：

由于且$X,Y$相互独立，所以$X,Y^2$相互独立、$X^2,Y$相互独立，于是：
\begin{align*}
E\left( XY^2-2X^2Y+1 \right) &=E\left( XY^2 \right) -2E\left( X^2Y \right) +E\left( 1 \right) \\
&=E\left( X \right) \cdot E\left( Y^2 \right) -2E\left( X^2 \right) \cdot E\left( Y \right) +1
\end{align*}
先分别计算$E\left( X \right) $、$E\left( X^2 \right) $、$E\left( Y \right) $和$E\left( Y^2 \right) $，然后计算结果：
\begin{align*}
&\because \begin{cases}
	E\left( X \right) =0\cdot \frac{3}{5}+10\cdot \frac{2}{5}=4\\
	E\left( X^2 \right) =0^2\cdot \frac{3}{5}+10^2\cdot \frac{2}{5}=40\\
	E\left( Y \right) =\int_{-\infty}^{+\infty}{y\frac{1}{2}e^{-\left| y \right|}\cdot dy}=0\\
	E\left( Y^2 \right) =\int_{-\infty}^{+\infty}{y^2\frac{1}{2}e^{-\left| y \right|}\cdot dy}=\int_0^{+\infty}{y^2e^{-y}\cdot dy}=2\\
\end{cases} \\
&\therefore E\left( X \right) \cdot E\left( Y^2 \right) -2E\left( X^2 \right) \cdot E\left( Y \right) +1=4\cdot 2-2\cdot 40\cdot 0+1=9
\end{align*}






\newpage
\section{方差}

本节介绍随机变量的第二个重要数字特征——方差。

本节要点：
\begin{itemize}
    \item 掌握方差的概念；
    \item 理解方差的统计意义和物理意义。
\end{itemize}

%============================================================
\subsection{一维随机变量方差的概念}

\begin{definition}[一维随机变量的方差]
设一维随机变量$X$，若函数$\left( X-EX \right) ^2$的数学期望存在，则称该数学期望为{\bf 一维随机变量$X$的方差}，记作$D\left( X \right) $或$DX$，即：
\begin{align*}
&DX:=E\left[ \left( X-EX \right) ^2 \right] =\sum_{i=1}^{+\infty}{\left( x_i-EX \right) ^2\cdot p_i} \\
&DX:=E\left[ \left( X-EX \right) ^2 \right] =\int_{-\infty}^{+\infty}{\left( x-EX \right) ^2f\left( x \right) \cdot dx}
\end{align*}
同时称$\sqrt{DX}$为{\bf $X$的平方差}。
\end{definition}

\begin{tcolorbox}
方差的英文为Variance，“教材\cite{book2}”中使用$\mathrm{Var}\left( X \right) $作为记号。
\end{tcolorbox}

方差的求解较为繁琐，由于$\left( X-EX \right) ^2=X^2-2\cdot X\cdot EX+\left( EX \right) ^2$，且$EX$是常数，利用数学期望的性质可得便捷公式：
\begin{align*}
DX&=E\left[ X^2-2\cdot X\cdot EX+\left( EX \right) ^2 \right] \\
&=E\left( X^2 \right) -2\cdot EX\cdot EX+\left( EX \right) ^2 \\
&=E\left( X^2 \right) -\left( EX \right) ^2
\end{align*}

从便捷公式可以得到$\left[ E\left( X \right) \right] ^2$和$E\left( X^2 \right) $的关系：
\[
E\left( X^2 \right) =DX+\left( EX \right) ^2
\]
由于$DX\geqslant 0$，所以$E\left( X^2 \right) \geqslant \left( EX \right) ^2$。

\begin{tcolorbox}
由于数学期望的线性性，便捷公式的推导没有任何前置条件，所以被广泛使用，定义式则更多用来理解概念。
\end{tcolorbox}

%============================================================
\subsection{二维随机变量函数的方差}

\begin{tcolorbox}
在二维变量中，我们更多考虑两个变量组合函数的方差。
\end{tcolorbox}

\begin{definition}[二维随机变量函数的方差]
设二维随机变量$\left( X,Y \right) $，若有$Z=z\left( X,Y \right) $且根据定义的对应级数或积分绝对收敛，具体定义式略，则称相应的级数或积分为{\bf $Z=z\left( X,Y \right) $的方差}，可采用便利公式计算，如下：
\[
DZ=E\left[ z\left( X,Y \right) ^2 \right] -E\left[ z\left( X,Y \right) \right] ^2
\]
\end{definition}

%============================================================
\subsection{方差的数学由来和统计意义}

方差定义中，$X-EX$是根本，表示了随机变量$X$和它的均值$EX$的偏离程度。
由于$E\left( X-EX \right) =E\left( X \right) -EX=0$，总是为0，所以我们取绝对值$\left| X-EX \right|$，又由于绝对值函数在0点不可导，所以我们替换成平方$\left( X-EX \right) ^2$定义方差。

方差描述了随机变量$X$偏离其数学期望的程度，$DX$越小，说明离散程度越小。

%============================================================
\subsection{方差的物理意义}

考虑一维物体的转动。
长度为$L$的细长杆质量$M$且分布不均匀，有线密度$\rho \left( x \right) ,x\in \left[ d,d+L \right] $，一端距原点$d$，绕原点转动，如下图。
\begin{figure}[h]
\centering
\includegraphics[height=1.5cm]{4.6.png}
\end{figure}

其转动惯量$I$有：
\[
I=\int{x^2\cdot dM}=\int_d^{d+L}{x^2\cdot \rho \left( x \right) \cdot dx} \quad \left[ \mathrm{kg}\cdot \mathrm{m}^2 \right]
\]
若将转轴移到质心$x_0$ ，根据平行轴定理，新的转动惯量$I_{x_0}$有：
\begin{align*}
&\because I=I_{x_0}+{Mx_0}^2 \\
&\therefore \int_d^{d+L}{x^2\cdot \rho \left( x \right) \cdot dx}=I_{x_0}+{Mx_0}^2 \\
&\therefore \int_d^{d+L}{x^2\cdot \frac{\rho \left( x \right)}{M}\cdot dx}=\frac{I_{x_0}}{M}+{x_0}^2
\end{align*}
对比方差的便捷计算公式：
\[
E\left( X^2 \right) =\int_{-\infty}^{+\infty}{x^2f\left( x \right) \cdot dx}=DX+\left( EX \right) ^2
\]
显然，{\bf 方差的物理意义可以理解为物体以质心为原点时的单位质量的转动惯量，有量纲$\left[ \mathrm{m}^2 \right] $}。
跟物理的转动惯量不一样，没有质量这个因子，所以描述的是一种和具体质量无关的惰性，这种惰性和分布有关，$DX$越大，说明分布越靠外，转动惯量也越大。

考虑二维平面，面密度$\rho \left( x,y \right) $的平面$D$，绕{\it y}轴匀速转动，如下左图：
\begin{figure}[h]
\centering
\includegraphics[height=3cm]{4.7.png}
\end{figure}

其转动惯量$I$有：
\[
I=\iint\limits_D{x^2\cdot \rho \left( x,y \right) \cdot d\sigma} \quad \left[ \mathrm{kg}\cdot \mathrm{m}^2 \right]
\]
对比二维变量的方差：
\[
DX=\iint\limits_{\mathbb{R} \times \mathbb{R}}{\left( x-EX \right) ^2f\left( x,y \right) \cdot d\sigma}
\]
{\bf 二维变量$\left( X,Y \right) $的方差$DX$的物理意义为在绕轴$x=x_0$竖直转动时，单位质量的转动惯量，有量纲$\left[ \mathrm{m}^2 \right] $。}
同样方差$DY$的物理意义为在绕轴$y=y_0$水平转动时，单位质量的转动惯量，也有量纲$\left[ \mathrm{m}^2 \right] $，如上右图。

%============================================================
\subsection{方差的性质}

{\bf 性质1}：常数的方差为0，即$D\left( c \right) =0$。

{\bf 推论}：$D\left( X \right) \geqslant 0$，当且仅当$P\left\{ X=EX \right\} =1$时（即常数分布），等号成立。

{\bf 性质2}：$D\left( kX \right) =k^2\cdot DX$

{\bf 性质3}：$D\left( kX+c \right) =k^2\cdot DX$，由性质1，2可得。

{\bf 性质4}：设$\left( X,Y \right) $为二维随机变量，

若$X,Y$相互独立$\Rightarrow D\left( X\pm Y \right) =DX+DY$

进一步地，若$X,Y$相互独立$\Rightarrow D\left( aX\pm bY \right) =a^2\cdot DX+b^2\cdot DY$

~

性质1，对于常数分布，质心就是$c$，显然绕质心的转动惯量为0。

性质2，若$X\rightarrow kX$，相当于直线被拉长$k$倍，自然转动惯量扩大了$k^2$倍。

\begin{figure}[h]
\centering
\includegraphics[height=2.5cm]{4.8.png}
\end{figure}

\begin{tcolorbox}
上述性质可结合物理意义理解。
\end{tcolorbox}

性质4，证明如下：
\begin{align*}
&\because DX=E\left( X^2 \right) -\left( EX \right) ^2 \\
&\begin{aligned}
	\therefore D\left( X\pm Y \right) &=E\left[ \left( X\pm Y \right) ^2 \right] -E\left[ \left( X\pm Y \right) \right] \cdot E\left[ \left( X\pm Y \right) \right]\\
	&=E\left[ X^2\pm 2XY+Y^2 \right] -\left[ EX\pm EY \right] ^2\\
	&=\left[ E\left( X^2 \right) \pm 2E\left( XY \right) +E\left( Y^2 \right) \right] -\\
	&\left[ \left( EX \right) ^2\pm 2\left( EX \right) \left( EY \right) +\left( EY \right) ^2 \right]\\
\end{aligned} \\
&\because E\left( XY \right) =\left( EX \right) \left( EY \right) \\
&\begin{aligned}
	\therefore D\left( X\pm Y \right) &=\left[ E\left( X^2 \right) -\left( EX \right) ^2 \right] +\left[ E\left( Y^2 \right) -\left( EY \right) ^2 \right]\\
	&=DX+DY\\
\end{aligned}
\end{align*}

%============================================================
\subsection{方差的取值范围}

变量的方差的取值范围根据不同的分布有不同的范围。如常数分布，中心点转动惯量必然为0。
又如$P\left\{ X=\pm \infty \right\} =0.5$，显然转动惯量$+\infty $ 。

\begin{figure}[h]
\centering
\includegraphics[height=2.5cm]{4.9.png}
\end{figure}

%============================================================
\subsection{随机变量的标准化}

\begin{definition}[随机变量的标准化]
设随机变量$X$有数学期望$EX$和方差$DX>0$，若将$X$进行运算$\frac{X-EX}{\sqrt{DX}}$，获得的随机变量称为{\bf $X$的标准化随机变量}，记作$X^*$，即：
\[
X^*=\frac{X-EX}{\sqrt{DX}}
\]
且有：
\begin{align*}
&E\left( X^* \right) =E\left( \frac{X-EX}{\sqrt{DX}} \right) =\frac{1}{\sqrt{DX}}\cdot E\left( X-EX \right) =0 \\
&D\left( X^* \right) =D\left( \frac{X-EX}{\sqrt{DX}} \right) =\frac{1}{DX}\cdot D\left( X-EX \right) =1
\end{align*}
\end{definition}

~

考察标准化过程：
\begin{itemize}
    \item $X-EX$:表示将坐标原点拉到期望$EX$的位置，即原点对称；
    \item $\frac{1}{\sqrt{DX}}$：拉伸或收缩{\it x}轴，使得分布不过于稀疏或稠密。
\end{itemize}
排除这两个“干扰”后，很多初看起来不同的分布，实际是一类分布，有利于后续的进一步分析。
如不同的正态分布，标准化之后就是标准正态分布。






\newpage
\section{常见分布的数学期望和方差}

本节讨论几个常见分布的数学期望和方差。

%============================================================
\subsection{0-1分布}

0-1分布$B\left( 1,p \right) $，$P\left\{ X=x \right\} =p^x\left( 1-p \right) ^{1-x},\quad x=0,1$：
\begin{align*}
&EX=p \\
&DX=p-p^2
\end{align*}

%============================================================
\subsection{二项分布}

二项分布$B\left( n,p \right) $，$P\left\{ X=x \right\} =C_{n}^{x}p^x\left( 1-p \right) ^{n-x},\quad x=0,1,\cdots ,n$：
\begin{align*}
&EX=np \\
&DX=np\left( 1-p \right)
\end{align*}

%============================================================
\subsection{泊松分布}

泊松分布$P\left( \lambda \right) $，$P\left\{ X=x \right\} =\frac{\lambda ^x}{x!}e^{-\lambda},\quad x=0,1,2,\cdots $：
\begin{align*}
&EX=\lambda \\
&DX=\lambda
\end{align*}
配合例图，$\lambda =np$越大，质心越远，转动惯量也越大。
\begin{figure}[h]
\centering
\includegraphics[height=5cm]{2.5.png}
\end{figure}

%============================================================
\subsection{几何分布}

几何分布$G\left( p \right) $，$P\left\{ X=x \right\} =p\left( 1-p \right) ^{x-1},\quad x=1,2,\cdots $：
\begin{align*}
&EX=\frac{1}{p} \\
&DX=\frac{1-p}{p^2}
\end{align*}
配合例图，$p$越大，质心越近，转动惯量也越小。
\begin{figure}[h]
\centering
\includegraphics[height=5cm]{2.6.png}
\end{figure}

%============================================================
\subsection{超几何分布}

超几何分布$H\left( M,N,n \right) $，$P\left\{ X=x \right\} =\frac{C_{N-M}^{n-x}C_{M}^{x}}{C_{N}^{n}}$：
\begin{align*}
&EX=\frac{nM}{N} \\
&DX=\frac{nM\left( N-n \right) \left( N-M \right)}{N^2\left( N-1 \right)}
\end{align*}

%============================================================
\subsection{均匀分布}
均匀分布$U\left[ a,b \right] $，$f\left( x \right) =\begin{cases}
	\frac{1}{b-a},&		x\in \left[ a,b \right]\\
	0,&		x\notin \left[ a,b \right]\\
\end{cases}$：
\begin{align*}
&EX=\frac{a+b}{2} \\
&DX=\frac{1}{12}\left( b-a \right) ^2
\end{align*}

%============================================================
\subsection{指数分布}

指数分布$E\left( \lambda \right) $，$f\left( x \right) =\begin{cases}
	0,&		x<0\\
	\lambda e^{-\lambda x},&		x\geqslant 0\\
\end{cases}$：
\begin{align*}
&EX=\frac{1}{\lambda} \\
&DX=\frac{1}{\lambda ^2}
\end{align*}
配合例图，$\lambda $越大，质心越近，转动惯量也越小。
\begin{figure}[h]
\centering
\includegraphics[height=5cm]{2.9.png}
\end{figure}

考察标准化指数分布：
\begin{align*}
&\because \begin{cases}
	X^*=\frac{X-EX}{\sqrt{DX}}\\
	X=\sqrt{DX}\cdot X^*+EX\\
\end{cases} \\
&\therefore \begin{cases}
	g^{-1}\left( X^* \right) =\sqrt{DX}\cdot X^*+EX\\
	\left| \frac{dg^{-1}\left( Y \right)}{dY} \right|=\sqrt{DX}\\
\end{cases} \\
&\begin{aligned}
	\therefore f\left( x^* \right) &=f_X\left[ g^{-1}\left( X^* \right) \right] \cdot \left| \frac{dg^{-1}\left( X^* \right)}{dX^*} \right|\\
	&=\lambda e^{-\lambda \left( \sqrt{DX}\cdot x^*+EX \right)}\cdot \sqrt{DX}=\lambda e^{-\lambda \left( \frac{1}{\lambda}\cdot x^*+\frac{1}{\lambda} \right)}\cdot \frac{1}{\lambda}\\
	&=e^{-\left( x^*+1 \right)}\\
\end{aligned}
\end{align*}
其中，$x^*\geqslant -1$，所以标准化的分布函数：
\[
F\left( x^* \right) =\int_{-1}^{+\infty}{f\left( t \right) \cdot dt}=\int_{-1}^{x^*}{e^{-\left( t+1 \right)}\cdot dt}=1-e^{-\left( x^*+1 \right)}
\]
\begin{figure}[ht]
\centering
\includegraphics[height=2.5cm]{4.10.png}
\end{figure}

%============================================================
\subsection{正态分布}

正态分布$N\left( \mu ,\sigma ^2 \right) $，$f\left( x \right) =\frac{1}{\sqrt{2\pi}\sigma}e^{-\frac{\left( x-\mu \right) ^2}{2\sigma ^2}} \quad \mu \in \mathbb{R} ,\sigma >0$：
\begin{align*}
&EX=\mu \\
&DX=\sigma ^2
\end{align*}
标准正态分布$N\left( 0,1 \right) $，$\varphi \left( x \right) =\frac{1}{\sqrt{2\pi}}e^{-\frac{x^2}{2}}$：
\begin{align*}
&EX=0 \\
&DX=1
\end{align*}
配合例图，$\mu $决定了质心的位置，$\sigma ^2$决定了转动惯量。

\begin{figure}[h]
\centering
\includegraphics[height=7.5cm]{2.10.png}
\end{figure}






\newpage
\section{协方差和相关系数}

前面介绍的数学期望和方差，无论是否为二维变量，最终考察的总是一维变量的数字特征。
本节介绍描述两个变量之间关系的数字特征——协方差和相关系数。

本节要点：
\begin{itemize}
    \item 掌握协方差的概念；
    \item 掌握相关系数的概念；
    \item 深刻理解相关系数的由来及几何意义。
\end{itemize}

%============================================================
\subsection{协方差的概念}

\begin{definition}[协方差]
设二维随机变量$\left( X,Y \right) $，若存在数学期望
\[
E\left[ \left( X-EX \right) \left( Y-EY \right) \right]
\]
则称之为{\bf $X$和$Y$的协方差}，记作$\mathrm{Cov}\left( X,Y \right) $，即：
\begin{align*}
\mathrm{Cov}\left( X,Y \right) :&=\sum_i{\sum_j{\left( x_i-EX \right) \cdot \left( y_j-EY \right) \cdot p_{ij}}} \\
\mathrm{Cov}\left( X,Y \right) :&=\iint\limits_{\mathbb{R} \times \mathbb{R}}{\left( x-EX \right) \cdot \left( y-EY \right) \cdot f\left( x,y \right) \cdot d\sigma}, \quad d\sigma =dxdy
\end{align*}
\end{definition}

由数学期望和方差的性质可得便捷计算公式：
\begin{align*}
\mathrm{Cov}\left( X,Y \right) &=E\left[ \left( X-EX \right) \left( Y-EY \right) \right] \\
&=E\left[ XY-X\cdot EY-Y\cdot EX+EX\cdot EY \right] \\
&=E\left( XY \right) -EY\cdot E\left( X \right) -EX\cdot E\left( Y \right) +E\left( EX\cdot EY \right) \\
&=E\left( XY \right) -EX\cdot EY
\end{align*}
特别地，当$X$和$Y$相互独立时$E\left( XY \right) =E\left( X \right) \cdot E\left( Y \right) $，有：
\[
\mathrm{Cov}\left( X,Y \right) =E\left( X \right) \cdot E\left( Y \right) -EX\cdot EY=0
\]

\begin{tcolorbox}
$\mathrm{Cov}$是“Covariance”的缩写。同样定义式更多的是为了理解概念，实际中通常使用便捷公式。
\end{tcolorbox}

%============================================================
\subsection{协方差的性质}

{\bf 性质1}：$\mathrm{Cov}\left( X,X \right) =DX$

{\bf 性质2}：$\mathrm{Cov}\left( X,c \right) =0$

{\bf 性质3}：设$\left( X,Y \right) $为二维随机变量，

$D\left( X\pm Y \right) =DX+DY\pm 2\mathrm{Cov}\left( X,Y \right) $

{\bf 性质4}：设$\left( X,Y \right) $为二维随机变量，$\mathrm{Cov}\left( X,Y \right) =\mathrm{Cov}\left( Y,X \right) $

{\bf 性质5}：设$\left( X,Y \right) $为二维随机变量，$\left[ \mathrm{Cov}\left( X,Y \right) \right] ^2\leqslant DX\cdot DY$

{\bf 性质6}：设$\left( X,Y \right) $为二维随机变量，$\mathrm{Cov}\left( aX,bY \right) =ab\mathrm{Cov}\left( X,Y \right) $

{\bf 性质7}：设$\left( X_1,X_2,Y \right) $为三维随机变量，

$\mathrm{Cov}\left( X_1\pm X_2,Y \right) =\mathrm{Cov}\left( X_1,Y \right) \pm \mathrm{Cov}\left( X_2,Y \right) $

~

性质1，用各自的便捷公式即可证明：
\[
\mathrm{Cov}\left( X,X \right) =E\left( X^2 \right) -\left( EX \right) ^2=DX
\]

性质3，用各自的便捷公式也可证明：
\begin{align*}
D\left( X\pm Y \right) &=\left[ E\left( X^2 \right) \pm 2E\left( XY \right) +E\left( Y^2 \right) \right] - \\
&\left[ \left( EX \right) ^2\pm 2\left( EX \right) \left( EY \right) +\left( EY \right) ^2 \right] \\
&=DX+DY\pm 2\left[ E\left( XY \right) +\left( EX \right) \left( EY \right) \right] \\
&=DX+DY\pm 2\mathrm{Cov}\left( X,Y \right)
\end{align*}

\begin{tcolorbox}
上述两个性质联系了方差和协方差，特别是性质3，可配合方差性质：
\[
\text{若}X,Y\text{相互独立}\Rightarrow D\left( X\pm Y \right) =D\left( X \right) +D\left( Y \right)
\]
进行理解。
\end{tcolorbox}

性质5，是理解相关系数的关键，证明如下：
\begin{align*}
&\because \mathrm{Cov}\left( X,Y \right) =E\left[ \left( X-EX \right) \left( Y-EY \right) \right] \\
&\because \left[ E\left( XY \right) \right] ^2\leqslant E\left( X^2 \right) \cdot E\left( Y^2 \right) \\
&\begin{aligned}
	\therefore \left[ \mathrm{Cov}\left( X,Y \right) \right] ^2&=\left[ E\left[ \left( X-EX \right) \left( Y-EY \right) \right] \right] ^2\\
	&\leqslant E\left[ \left( X-EX \right) ^2 \right] \cdot E\left[ \left( Y-EY \right) ^2 \right] =DX\cdot DY\\
\end{aligned}
\end{align*}

%============================================================
\subsection{协方差的统计意义}
从协方差的定义入手：
\[
\mathrm{Cov}\left( X,Y \right) :=\sum_i{\sum_j{\left( x_i-EX \right) \cdot \left( y_j-EY \right) \cdot p_{ij}}}
\]

\begin{itemize}
    \item $\left( x_i-EX \right) \cdot \left( y_j-EY \right) $：表示样本点对均值构成的矩形面积，且1、3象限取正，2、4象限取负；
    \item $\sum_i{\sum_j{\cdot p_{ij}}}$：表示对这些面积加权和。
\end{itemize}
可见：
\begin{itemize}
    \item 当$\mathrm{Cov}\left( X,Y \right) >0$时，表示1、3象限居多，统计上，对于一个样本，$X$取大值时，可能$Y$也是大值；
    \item 反之，当$\mathrm{Cov}\left( X,Y \right) <0$时，统计上，对于一个样本，$X$取大值时，$Y$反倒可能是小值。
\end{itemize}
协方差就表示了$X$和$Y$在分布上有的这个关系，但协方差没有去除$X,Y$各自的方差，下面的相关系数除去了方差的影响，更能精准地刻画$X$和$Y$在分布上的关系。

%============================================================
\subsection{相关系数和不相关的概念}

\begin{definition}[相关系数]
设二维随机变量$\left( X,Y \right) $，若$DX>0,DY>0$，则称$\frac{\mathrm{Cov}\left( X,Y \right)}{\sqrt{DX}\cdot \sqrt{DY}}$为{\bf 二维随机变量$\left( X,Y \right) $的相关系数}，记作$\rho _{XY}$或$\rho $，即：
\[
\rho _{XY}:=\frac{\mathrm{Cov}\left( X,Y \right)}{\sqrt{DX}\cdot \sqrt{DY}}
\]
\end{definition}

由各自的便捷公式可得相关系数的便捷公式：
\[
\rho _{XY}=\frac{\mathrm{Cov}\left( X,Y \right)}{\sqrt{DX}\cdot \sqrt{DY}}=\frac{E\left( XY \right) -EX\cdot EY}{\sqrt{E\left( X^2 \right) -\left( EX \right) ^2}\cdot \sqrt{E\left( Y^2 \right) -\left( EY \right) ^2}}
\]

\begin{definition}[不相关]
对于二维变量$\left( X,Y \right) $，若$\rho _{XY}=0$，则称它们{\bf 不相关}。
\end{definition}

不相关指的是$X,Y$没有线性关系，不排除有其他关系，如平方关系，三角函数关系等，所以不相关是一个比独立要弱的要求。
$X,Y$不相关并不一定独立，而独立必然不相关。

\begin{theorem}
若$X,Y$独立且$\rho _{XY}$存在，则$X,Y$不相关。
\end{theorem}

简单来说即，独立$\Rightarrow $不相关，还是要注意反命题不成立。

\begin{theorem}
若二维变量$\left( X,Y \right) $存在$\rho _{XY}$，则下列命题等价：
\begin{itemize}
    \item $X,Y$不相关；
    \item $\rho _{XY}=0$
    \item $\mathrm{Cov}\left( X,Y \right) =0$
    \item $E\left( XY \right) =EX\cdot EY$
    \item $D\left( X\pm Y \right) =D\left( X \right) +D\left( Y \right) $
\end{itemize}
\end{theorem}

%============================================================
\subsection{相关系数的性质}

{\bf 性质1}：$\rho _{XY}\in \left[ -1,1 \right] $

{\bf 性质2}：$\rho _{XY}=\pm 1$的充要条件是$X$和$Y$呈线性关系，即$\exists a\ne 0,b$使得$P\left\{ Y=aX+b \right\} =1$，且此时有$a>0\Leftrightarrow \rho _{XY}=1,a<0\Leftrightarrow \rho _{XY}=-1$。

{\bf 性质3}：$\forall a\ne 0,\forall b\ne 0$，有
\[
\rho _{\left( aX \right) \left( bY \right)}=\begin{cases}
	\rho _{XY},&		ab>0\\
	-\rho _{XY},&		ab<0\\
\end{cases}
\]

性质1，由相关系数的定义和协方差的性质7可证。

性质2在下面说明。

%============================================================
\subsection{相关系数的由来}

以二维离散变量$\left( X,Y \right) $为例说明，首先我们定义一个函数：
\[
Z=z\left( X,Y \right) =\left[ Y-\left( aX+b \right) \right] ^2
\]
我们求它的数学期望：
\[
EZ=E\left[ z\left( X,Y \right) \right] =\sum_i{\sum_j{\left[ Y-\left( aX+b \right) \right] ^2\cdot p_{ij}}}
\]
若$\left( X,Y \right) $的分布如下左图，我们用圆圈表示$p_{ij}$，圆圈越大表示$p_{ij}$越大，一共7个 ，呈一直线排列。
可见，只有当$X,Y$取值满足$y=ax+b$时，有$p_{ij}\ne 0$，于是$EZ=0$。
若$\left( X,Y \right) $的分布如下中图，4个不在一条直线上的点使得$EZ\ne 0$。
同样，如下右图，$EZ\ne 0$，而且$EZ$会变得更大。

\begin{figure}[h]
\centering
\includegraphics[height=2.2cm]{4.11.png}
\end{figure}

相关系数考察的就是这个关系。

首先引入线性代数的理论，令 $\hat{Y}=aX+b$，即$Y$在$X$上的投影：
\begin{align*}
e\left( a,b \right) &=E\left\{ \left[ Y-\left( aX+b \right) \right] ^2 \right\} \\
&=E\left( Y^2 \right) +a^2E\left( X^2 \right) +b^2-2aE\left( XY \right) -2bE\left( Y \right) +2abE\left( X \right)
\end{align*}

然后我们考察$X$需要做出怎么样的调整使得和$Y$的这个差异最小，即考察$e\left( a,b \right) $的极值。首先用一阶偏导求极值点$\left( a_0,b_0 \right) $：
\begin{align*}
&\begin{cases}
	\frac{\partial e}{\partial a}=2aE\left( X^2 \right) -2E\left( XY \right) +2bE\left( X \right) =0\\
	\frac{\partial e}{\partial b}=2b-2E\left( Y \right) +2aE\left( X \right) =0\\
\end{cases} \\
&\Rightarrow \begin{cases}
	a_0=\frac{\mathrm{Cov}\left( X,Y \right)}{D\left( X \right)}\\
	b_0=E\left( Y \right) -\frac{\mathrm{Cov}\left( X,Y \right)}{D\left( X \right)}\cdot E\left( X \right)\\
\end{cases}
\end{align*}
但这只是极值点，极大值还是极小值还是鞍点，需要考虑$e\left( a,b \right) $的Hessian矩阵的特征值：
\begin{align*}
&H\left( e \right) =\left[ \begin{matrix}
	\frac{\partial ^2e}{\partial a^2}&		\frac{\partial ^2e}{\partial a\partial b}\\
	\frac{\partial ^2e}{\partial b\partial a}&		\frac{\partial ^2e}{\partial b^2}\\
\end{matrix} \right] =\left[ \begin{matrix}
	2E\left( X^2 \right)&		2E\left( X \right)\\
	2E\left( X \right)&		2\\
\end{matrix} \right] \\
&\det \left[ H\left( e \right) \right] =4E\left( X^2 \right) -4\left[ 2E\left( X \right) \right] ^2=4D\left( X \right) >0
\end{align*}
再加上$\frac{\partial ^2e}{\partial a^2}=2E\left( X^2 \right) >0$，所以$\left( a_0,b_0 \right) $是极小值点，$e\left( a,b \right) $在该点有最小值：
\[
\left. e\left( a,b \right) \right|_{\left( a_0,b_0 \right)}=\left[ 1-\left( \frac{\mathrm{Cov}\left( X,Y \right)}{\sqrt{D\left( X \right)}\cdot \sqrt{D\left( Y \right)}} \right) ^2 \right] D\left( Y \right)
\]

上式衡量了$X$和$Y$在分布上的差异量，即 “努力地做出调整”后，$Y$和$X$的差异量。
若要考察一个已知的$Y$对$X_1$和$X_2$的线性关系谁大谁小，则需要排除$Y$的因素，于是我们可取前半部分因子
\[
\left[ 1-\left( \frac{\mathrm{Cov}\left( X_?,Y \right)}{\sqrt{D\left( X_? \right)}\cdot \sqrt{D\left( Y \right)}} \right) ^2 \right]
\]
该因子越小，说明差异量越小，和$X_?$的契合度越高。反之取平方项部分，记作$\rho _{XY}$
\[
\rho _{X_?Y}=\frac{\mathrm{Cov}\left( X_?,Y \right)}{\sqrt{D\left( X_? \right)}\cdot \sqrt{D\left( Y \right)}}
\]
$\rho _{X_?Y}$越大，说明和$X_?$的契合度越高。
也就是说，对于一个给定的$Y$和一堆$X_1,X_2,\cdots X_n$，我们可以用$\rho _{XY}$判断$Y$和某个$X_i$在分布上是最契合的，此时，$X_i$只需要做出调整$a_0X_i+b_0$就可以将差异减小到最小的程度。

对于现实的指导是，如果$\rho _{XY}$大到一定程度，我们就可以用关于$X$的一个简单的函数
\[
g\left( X \right) =\frac{\mathrm{Cov}\left( X,Y \right)}{D\left( X \right)}X+E\left( Y \right) -\frac{\mathrm{Cov}\left( X,Y \right)}{D\left( X \right)}\cdot E\left( X \right)
\]
代替$Y$考察分布。
比如对于正态分布，有很多很好很简单的性质，我们也有很多统计手段和公式。
倘若 “很像”正态分布，那我们就可以用正态分布代替$Y$，进而可以通过公式用标准分布计算，整个问题就变得简单明了。

%============================================================
\subsection{相关的几何意义}

$\rho _{XY}$表示的相关性如下图。
可见，当$\rho _{XY}=\pm 1$时，$Y\ne aX+b$的地方概率为0，$f\left( x,y \right) $将由一张三维曲面退化成一个二维平面。

\begin{figure}[h]
\centering
\includegraphics[height=4.5cm]{4.12.png}
\end{figure}

以二维正态分布为例：
\begin{align*}
&\left( X,Y \right) \sim N\left( \mu _1,\mu _2,{\sigma _1}^2,{\sigma _2}^2,\rho \right) \\
&f\left( x,y \right) =\frac{1}{2\pi \sigma _1\sigma _2\sqrt{1-\rho ^2}}e^{-\frac{1}{2\left( 1-\rho ^2 \right)}\left[ \frac{\left( x-\mu _1 \right) ^2}{{\sigma _1}^2}-2\rho \frac{\left( x-\mu _1 \right) \left( y-\mu _2 \right)}{\sigma _1\sigma _2}+\frac{\left( y-\mu _2 \right) ^2}{{\sigma _2}^2} \right]} \\
&\rho \in \left( -1,1 \right)
\end{align*}
$\rho $表示了$X,Y$的相关性，当$\rho =0.99$时，$f\left( x,y \right) $几乎成为一个平面。

\begin{figure}[h]
\centering
\includegraphics[height=8cm]{4.13.png}
\end{figure}

\begin{tcolorbox}
可见，相关描述了密度函数的扁平性，相关系数就是这个扁平性的量化指标，$\rho _{XY}$越大，曲面$f\left( x,y \right) $越趋于平面。
\end{tcolorbox}

%============================================================
\subsection{协方差和相关系数的线性代数意义}

参考线性代数中向量的柯西—施瓦茨不等式的表述。若$\boldsymbol{x},\boldsymbol{y}$为内积空间$V$中的两个向量，则必有：
\[
\left| \left< \boldsymbol{x},\boldsymbol{y} \right> \right|\leqslant \sqrt{\left< \boldsymbol{x},\boldsymbol{x} \right>}\sqrt{\left< \boldsymbol{y},\boldsymbol{y} \right>}
\]
当且仅当$\boldsymbol{x}=0$或$\boldsymbol{y}=0$或$\boldsymbol{x}=\boldsymbol{y}=0$或$\boldsymbol{x},\boldsymbol{y}$线性相关时等号成立。
于是当$\boldsymbol{x}\ne \boldsymbol{y}$时，我们可以用柯西—施瓦茨不等式定义它们的夹角：
\[
\cos \theta =\frac{\left< \boldsymbol{x},\boldsymbol{y} \right>}{\sqrt{\left< \boldsymbol{x},\boldsymbol{x} \right>}\sqrt{\left< \boldsymbol{y},\boldsymbol{y} \right>}}
\]
对比观察得到，协方差就是两个向量的内积，相关系数就是它们的夹角的余弦：
\[
\rho _{XY}=\frac{\mathrm{Cov}\left( X,Y \right)}{\sqrt{DX}\cdot \sqrt{DY}}=\frac{\mathrm{Cov}\left( X,Y \right)}{\sqrt{\mathrm{Cov}\left( X,X \right)}\cdot \sqrt{\mathrm{Cov}\left( Y,Y \right)}}
\]
各部分对应关系如下：
\[
\begin{array}{l}
	\mathrm{Cov}\left( , \right)\\
	X\\
	Y\\
	DX\\
	DY\\
	\mathrm{Cov}\left( X,Y \right)\\
	\rho _{XY}\\
\end{array} \quad \leftrightarrow \quad \begin{array}{l}
	\left< , \right>\\
	\boldsymbol{x}\\
	\boldsymbol{y}\\
	\left< \boldsymbol{x},\boldsymbol{x} \right>\\
	\left< \boldsymbol{y},\boldsymbol{y} \right>\\
	\left< \boldsymbol{x},\boldsymbol{y} \right>\\
	\cos \theta\\
\end{array}
\]

更进一步，将一维随机变量的分布视为向量，则所有分布的集合可以视作向量空间$V$，任取两种分布$X,Y$用向量记号记作$\boldsymbol{x},\boldsymbol{y}$，我们采用协方差作为操作，即$\left< , \right> =\mathrm{Cov}\left( , \right) $，必然满足内积的4个要求：
\begin{enumerate}[label=\Roman*.]
    \item 正定性，$\left< \boldsymbol{x},\boldsymbol{x} \right> \geqslant 0$，当且仅当$\boldsymbol{x} =0$（即常数分布）时等号成立；
    \item $\left< \alpha \boldsymbol{x},\boldsymbol{y} \right> =\alpha \left< \boldsymbol{x},\boldsymbol{y} \right> $；
    \item 交换律，$\left< \boldsymbol{x},\boldsymbol{y} \right> =\left< \boldsymbol{y},\boldsymbol{x} \right> $；
    \item 分配律，$\left< \boldsymbol{x}+\boldsymbol{y},\boldsymbol{z} \right> =\left< \boldsymbol{x},\boldsymbol{z} \right> +\left< \boldsymbol{y},\boldsymbol{z} \right> $。
\end{enumerate}
其中I、II、III这3条从数学期望的定义式就可得
IV证明如下：
\begin{align*}
&\mathrm{Cov}\left( \left( X+Y \right) ,Z \right) =E\left[ \left( X+Y-E\left( X+Y \right) \right) \left( Z-EZ \right) \right] \\
&=E\left[ \left( X+Y+EX+EY \right) \left( Z-EZ \right) \right] \\
&=E\left( XZ-XEZ+YZ-YEZ+ZEX-EXEZ+ZWY-EZEY \right) \\
&=E\left( XZ \right) +E\left( YZ \right) -E\left( EXEZ \right) -E\left( EZEY \right) \\
&=\mathrm{Cov}\left( X,Z \right) +\mathrm{Cov}\left( Y,Z \right)
\end{align*}

\begin{tcolorbox}
可见，一维随机变量连同协方差运算构成了内积空间，常数分布是零向量，$\rho _{XY}$表示了两个向量的夹角的余弦。
\end{tcolorbox}

%============================================================
\subsection{方差与范数}

更“线性代数化地说”，方差是变量的一个范数的算术平方$DX\leftrightarrow \left\| \boldsymbol{x} \right\| ^2$。范数的要求如下：
\begin{itemize}
    \item 正定性，$\left\| \boldsymbol{x} \right\| \geqslant 0$，当且仅当$\boldsymbol{x}=0$时，等号成立；
    \item $\left\| \alpha \boldsymbol{x} \right\| =\left| \alpha \right|\left\| \boldsymbol{x} \right\| $；
    \item 三角不等式，$\left\| \boldsymbol{x}+\boldsymbol{y} \right\| \leqslant \left\| \boldsymbol{x} \right\| +\left\| \boldsymbol{y} \right\| $。
\end{itemize}
对应方差的性质：
\begin{itemize}
    \item $D\left( X \right) \geqslant 0$，当且仅当$P\left\{ X=EX \right\} =1$时等号成立。
    \item $D\left( kX \right) =k^2\cdot D\left( X \right) $
    \item $D\left( X+Y \right) =DX+DY+2\mathrm{Cov}\left( X,Y \right) $
\end{itemize}

%============================================================
\subsection{例}

\begin{example}
继续第三章中的例子，若5个产品中有2个次品，每次抽取一个，取两次，记$X$表示第一次取到次品的个数，$Y$表示第二次取到次品的个数，讨论：
\begin{enumerate}
    \item 不放回抽取的情况下，
    \item 放回抽取的情况下，
\end{enumerate}
这两种情况下，$X,Y$的协方差。
\end{example}

解：

由题干可得，$X,Y$取值只可能是0或1。

1. 不放回抽取的情况下的分布律$p_{ij}$：
\[
\begin{matrix}
	&		0&		1&		p_{\cdot j}\\
	0&		\frac{3}{10}&		\frac{3}{10}&		\frac{3}{5}\\
	1&		\frac{3}{10}&		\frac{1}{10}&		\frac{2}{5}\\
	p_{i\cdot}&		\frac{3}{5}&		\frac{2}{5}&		1\\
\end{matrix}
\]
\begin{align*}
&\because \begin{cases}
	E\left( XY \right) =0\cdot 0\cdot \frac{3}{10}+0\cdot 1\cdot \frac{3}{10}+1\cdot 0\cdot \frac{3}{10}+1\cdot 1\cdot \frac{1}{10}=\frac{1}{10}\\
	E\left( X \right) =0\cdot \frac{3}{5}+1\cdot \frac{2}{5}=\frac{2}{5}\\
	E\left( Y \right) =0\cdot \frac{3}{5}+1\cdot \frac{2}{5}=\frac{2}{5}\\
\end{cases} \\
&\therefore \mathrm{Cov}\left( X,Y \right) =E\left( XY \right) -E\left( X \right) \cdot E\left( Y \right) =\frac{1}{10}-\frac{4}{25}=-\frac{3}{50}
\end{align*}

2. 放回抽取的情况下的分布律$p_{ij}$：
\[
\begin{matrix}
	&		0&		1&		p_{\cdot j}\\
	0&		\frac{9}{25}&		\frac{6}{25}&		\frac{3}{5}\\
	1&		\frac{6}{25}&		\frac{4}{25}&		\frac{2}{5}\\
	p_{i\cdot}&		\frac{3}{5}&		\frac{2}{5}&		1\\
\end{matrix}
\]
\begin{align*}
&\because \begin{cases}
	E\left( XY \right) =0\cdot 0\cdot \frac{9}{25}+0\cdot 1\cdot \frac{6}{25}+1\cdot 0\cdot \frac{6}{25}+1\cdot 1\cdot \frac{4}{25}=\frac{4}{25}\\
	E\left( X \right) =0\cdot \frac{3}{5}+1\cdot \frac{2}{5}=\frac{2}{5}\\
	E\left( Y \right) =0\cdot \frac{3}{5}+1\cdot \frac{2}{5}=\frac{2}{5}\\
\end{cases} \\
&\therefore \mathrm{Cov}\left( X,Y \right) =E\left( XY \right) -E\left( X \right) \cdot E\left( Y \right) =\frac{4}{25}-\frac{4}{25}=0
\end{align*}

\begin{tcolorbox}
从该题解读协方差的意义。在(2)放回的情况下，从事实角度出发，$X,Y$是独立的，既然独立就没有任何关系，包括线性关系，自然$\mathrm{Cov}\left( X,Y \right) =0$。
在(1)不放回的情况下，第一次抽中次品时，是不利于再次抽中次品这个事件的发生的。
\end{tcolorbox}






\newpage
\section{矩和协方差矩阵}

本节简单介绍矩和协方差矩阵的概念。

本节要点：
\begin{itemize}
    \item 了解矩的概念；
    \item 了解协方差矩阵的概念。
\end{itemize}

%============================================================
\subsection{矩的概念}

\begin{definition}[矩]
设一维随机变量$X$，若对于$k\in \mathbb{Z} ^+$，$E\left[ \left( X-c \right) ^k \right] $存在，则称其为{\bf $X$关于$c$点的$k$阶矩}。
特别地，当$c=0$时，称$E\left( X^k \right) $为{\bf $X$的$k$阶原点矩}，当$c=EX$时，称$E\left[ \left( X-EX \right) ^k \right] $为{\bf $X$的$k$阶中心矩}。
\end{definition}

显然：
\begin{itemize}
    \item $X$的1阶原点矩即为数学期望，$EX$；
    \item $X$的1阶中心矩总是为0，$E\left( X-EX \right) =0$；
    \item $X$的2阶中心矩即为方差，$DX$。
\end{itemize}

\begin{definition}[混合矩]
设二维随机变量$\left( X,Y \right) $，若对于$k,l\in \mathbb{Z} ^+$，$E\left( X^kY^l \right) $存在，则称其为{\bf $X,Y$的$k+l$阶混合原点矩}；
若对于$k,l\in \mathbb{Z} ^+$，数学期望$E\left[ \left( X-EX \right) ^k\left( Y-EY \right) ^l \right] $存在，则称其为{\bf $X,Y$的$k+l$阶混合中心矩}。
\end{definition}

显然：
\begin{itemize}
    \item $X,Y$的1+1阶混合中心矩即为它们的协方差，$\mathrm{Cov}\left( X,Y \right) $。
\end{itemize}

\begin{tcolorbox}
1阶矩和2阶矩用得较多，3、4阶用的较少。
3、4阶能引出“偏度系数”和“峰度系数”，详见“教材\cite{book2}”P122～123。
广义上来讲，矩也是数字特征，但在实际使用中，数字特征往往特指期望、方差、协方差。
\end{tcolorbox}

%============================================================
\subsection{协方差矩阵的概念}

\begin{definition}[协方差矩阵]
设$n$维随机变量$\left( X_1,X_2,\cdots ,X_n \right) $，记它们随意两两组合的协方差为$c_{ij}$，即：
\[
c_{ij}:=\mathrm{Cov}\left( X_i,X_j \right)
\]
若$c_{ij}$总是存在，则称它们组成的矩阵为{\bf $\left( X_1,X_2,\cdots ,X_n \right) $的协方差矩阵}，记作$\mathbf{C}$，即：
\[
\mathbf{C}:=\left[ \begin{matrix}
	c_{11}&		c_{12}&		\cdots&		c_{1n}\\
   	c_{21}&		c_{22}&		\cdots&		c_{2n}\\
   	\vdots&		\vdots&		&		\vdots\\
   	c_{n1}&		c_{n2}&		\cdots&		c_{nn}\\
\end{matrix} \right]
\]
\end{definition}

显然$\mathbf{C}$是一个实对称方阵，其对角线为每个变量的方差。






\newpage
\section{本章小结}

本章介绍了随机变量的数字特征：
\begin{itemize}
    \item 对于一个变量，有：数学期望、方差，描述了变量的分布情况；
    \item 对于两个变量，有：协方差、相关系数，描述了样本点$\left\{ X=x_i,Y=y_j \right\} $的分布是否趋于直线。
\end{itemize}

~

数学期望定义公式
\begin{align*}
&EX:=\sum_i{x_i\cdot p_i} \\
&EX:=\int_{-\infty}^{+\infty}{xf\left( x \right) \cdot dx}
\end{align*}

方差定义公式和便捷公式
\begin{align*}
&DX:=E\left[ \left( X-EX \right) ^2 \right] \\
&DX=E\left( X^2 \right) -\left( EX \right) ^2
\end{align*}

协方差定义公式和便捷公式
\begin{align*}
&\mathrm{Cov}\left( X,Y \right) :=E\left[ \left( X-EX \right) \left( Y-EY \right) \right] \\
&\mathrm{Cov}\left( X,Y \right) =E\left( XY \right) -EX\cdot EY
\end{align*}

相关系数定义公式和便捷公式
\begin{align*}
&\rho _{XY}:=\frac{\mathrm{Cov}\left( X,Y \right)}{\sqrt{DX}\cdot \sqrt{DY}} \\
&\rho _{XY}=\frac{E\left( XY \right) -EX\cdot EY}{\sqrt{E\left( X^2 \right) -\left( EX \right) ^2}\cdot \sqrt{E\left( Y^2 \right) -\left( EY \right) ^2}}
\end{align*}

数学期望的性质
\begin{align*}
&E\left( c \right) =c \\
&X\geqslant a\Rightarrow EX\geqslant a,X\leqslant a\Rightarrow EX\leqslant a \\
&E\left( kX \right) =k\cdot E\left( X \right) \\
&E\left( kX+c \right) =k\cdot E\left( X \right) +c \\
&E\left( X\pm Y \right) =E\left( X \right) \pm E\left( Y \right) \\
&\text{若}X,Y\text{相互独立}\Rightarrow E\left( XY \right) =E\left( X \right) \cdot E\left( Y \right) \\
&\left[ E\left( XY \right) \right] ^2\leqslant E\left( X^2 \right) \cdot E\left( Y^2 \right)
\end{align*}

方差的性质
\begin{align*}
&D\left( c \right) =0 \\
&D\left( X \right) \geqslant 0, \text{当且仅当}P\left\{ X=EX \right\} =1\text{时等号成立} \\
&D\left( kX \right) =k^2\cdot D\left( X \right) \\
&D\left( kX+c \right) =k^2\cdot D\left( X \right) \\
&\text{若}X,Y\text{相互独立}\Rightarrow D\left( X\pm Y \right) =D\left( X \right) +D\left( Y \right)
\end{align*}

协方差的性质
\begin{align*}
&\mathrm{Cov}\left( X,X \right) =DX \\
&D\left( X\pm Y \right) =DX+DY\pm 2\mathrm{Cov}\left( X,Y \right) \\
&\mathrm{Cov}\left( X,c \right) =0 \\
&\mathrm{Cov}\left( X,Y \right) =\mathrm{Cov}\left( Y,X \right) \\
&\mathrm{Cov}\left( aX,bY \right) =ab\mathrm{Cov}\left( X,Y \right) \\
&\mathrm{Cov}\left( X_1\pm X_2,Y \right) =\mathrm{Cov}\left( X_1,Y \right) \pm \mathrm{Cov}\left( X_2,Y \right) \\
&\left[ \mathrm{Cov}\left( X,Y \right) \right] ^2\leqslant DX\cdot DY
\end{align*}

相关系数的性质
\begin{align*}
&\rho _{XY}\in \left[ -1,1 \right] \\
&X\text{和}Y\text{成线性关系}\Leftrightarrow \rho _{XY}=\pm 1 \\
&\rho _{\left( aX \right) \left( bY \right)}=\begin{cases}
	\rho _{XY},&		ab>0\\
	-\rho _{XY},&		ab<0\\
\end{cases}
\end{align*}

~

数字特征的物理意义和线性代数意义：
\begin{table}[h]
\centering
\begin{tabular}{lll}
    \toprule
    数字特征 & 物理意义/线性代数意义 & 量纲\\
    \midrule
    数学期望 & 质心位置 & $\left[ \mathrm{m} \right] $\\
    方差 & 单位质量的转动惯量，范数 & $\left[ \mathrm{m}^2 \right] $\\
    协方差 & 两个向量的内积 & $\left[ \mathrm{m}^2 \right] $\\
    相关系数 & 两个向量的夹角的余弦 & 无量纲\\
    \bottomrule
\end{tabular}
\end{table}

学习概率论中的数字特征，要特别注意它们的物理意义和几何意义，用其他学科的知识点相互印证，加深对这些概念的理解。









